%%%%%%%%%%%%%%%%%%%%%%%%%%%%%%%%%%%%%%%%%%%%%%%%%%%%%%%%%%%%%%%%%%%%%%%%
% Manuscript: Consciousness and the Cyclic Universe
% Author: Xiaohui Li
% License: CC-BY 4.0 for academic reuse
%%%%%%%%%%%%%%%%%%%%%%%%%%%%%%%%%%%%%%%%%%%%%%%%%%%%%%%%%%%%%%%%%%%%%%%%

\documentclass[11pt,a4paper]{article}

% Packages
\usepackage[utf8]{inputenc}
\usepackage[T1]{fontenc}
\usepackage{lmodern}
\usepackage{amsmath,amssymb,amsfonts}
\usepackage{graphicx}
\usepackage{booktabs}
\usepackage{natbib}
\usepackage[margin=2.5cm]{geometry}
\usepackage{hyperref}
\usepackage{xcolor}
\usepackage{lineno}
\usepackage{authblk}
\usepackage{titlesec}

% Citation style
\setcitestyle{numbers,square,comma}

% Line numbers for review
\linenumbers

% Title formatting
\titleformat{\section}{\normalfont\Large\bfseries}{\thesection}{1em}{}
\titleformat{\subsection}{\normalfont\large\bfseries}{\thesubsection}{1em}{}

% Title and author
\title{\textbf{Consciousness and the Cyclic Universe: A Quantum-Optical Hypothesis of Neural Information Processing and Cosmic Recurrence}}

\author[1]{Xiaohui Li}
\affil[1]{Independent Researcher, Baoji, China}
\affil[ ]{Correspondence: Xiaohui Li (email: \texttt{danceluck@163.com}); ORCID: 0009-0003-8233-7191}
\affil[ ]{Data and code homepage: \url{https://github.com/xiaohui-li-lab/}}

\date{}

\begin{document}

\maketitle

\begin{abstract}
\noindent
How consciousness emerges from the material brain and how it may be connected to the origin and fate of the cosmos remain among the deepest open questions in science. Here we propose the Infrared Quantum Information Channel (IRQIC) hypothesis: neural information processing is not limited to classical electrochemical signaling (``coarse quanta''), but also relies on a ``fine-quantum'' layer carried by near- to far-infrared photons. These photons are hypothesized to perform decomposition-recombination computations via quantum superposition inside subcellular structures such as neural microtubules. When consciousness is cultivated through meditation, contemplation, or inward stillness, neural noise is reduced and fine-quantum coherence may be enhanced, producing testable signatures of ``mind-as-quantum-state.'' We extend this microscopic model to cosmology: the Big Bang singularity is not an absolute information vacuum, but an information-bearing ``consciousness node'' embedded in a spacetime-transcending consciousness field. The universe evolves toward a ``Great Refinement,'' condenses into black-hole seeds, and a new universe is born from inside the horizon, forming a consciousness-driven cyclic cosmos. We provide key theoretical equations, numerical simulation protocols, falsifiable predictions, benchmark comparisons, and open-source code with synthetic datasets.
\end{abstract}

\noindent\textbf{Keywords:} consciousness; infrared photon; neural microtubule; quantum information; cyclic universe; big bounce; consciousness field

\section{Introduction}

Classical neuroscience has provided a detailed account of brain information transfer. Action potentials propagate along nerve fibers through voltage-gated Na$^+$ and K$^+$ channels: unmyelinated fibers conduct at 0.5--2 m/s, whereas myelinated fibers achieve saltatory conduction up to $\sim$120 m/s \cite{Hodgkin1952,Waxman1980}. At synapses, Ca$^{2+}$ influx triggers vesicle fusion and neurotransmitter release (e.g., acetylcholine, glutamate), producing excitatory or inhibitory postsynaptic potentials \cite{Katz1969}. These processes can be collectively described as ``coarse-quantum'' information processing, because their carriers---ions, water molecules, neurotransmitters, and cells---are well approximated by classical or semiclassical physics.

Yet the coarse-quantum model faces the ``hard problem'' of consciousness: why should electrochemical dynamics give rise to subjective experience \cite{Chalmers1995}? In addition, quantum coherence in microtubules, ultra-weak photon emission from biological systems, and Fr\"{o}hlich condensation suggest that a more subtle information layer may operate in living tissue \cite{Hameroff2014,Craddock2018,Popp1992,Frohlich1968}. Li has proposed that infrared photons constitute the core carrier of this ``fine-quantum'' layer \cite{Li2026a}. Here we formalize this proposal into an academic framework and explore its connection to cyclic cosmology.

\section{Theoretical Framework: A Dual-Scale Model of Neural Information Processing}

\subsection{Coarse-quantum layer}

The coarse-quantum layer corresponds to established mechanisms:

\paragraph{(i) Axonal electrical conduction.} Upon reaching threshold, voltage-gated Na$^+$ channels open, Na$^+$ rushes in and depolarizes the membrane; Na$^+$ inactivation and K$^+$ efflux then restore the resting potential \cite{Hodgkin1952,Waxman1980}.

\paragraph{(ii) Synaptic chemical transmission.} Arrival of the action potential at the axon terminal opens Ca$^{2+}$ channels; Ca$^{2+}$ triggers vesicle fusion, neurotransmitter release into the cleft, receptor binding on the postsynaptic membrane, and altered ion permeability \cite{Katz1969}.

\subsection{Fine-quantum layer: the Infrared Quantum Information Channel (IRQIC)}

The fine-quantum layer uses photons, especially infrared photons (wavelength $\sim$0.75--1000 $\mu$m, energy $\sim$0.001--1.65 eV), as information carriers. We ascribe three functions to them:

\paragraph{(i) Penetrative information transfer.} Infrared photons can propagate through the tubulin wall and surrounding water layers, forming low-loss waveguide modes inside microtubules \cite{Pokorny2004}. With an outer diameter of $\sim$25 nm and inner diameter of $\sim$15 nm, a microtubule is comparable in scale to infrared wavelengths and may support dielectric waveguide modes \cite{Craddock2018,Pokorny2004}.

\paragraph{(ii) Superposition-state computation.} Information is encoded in infrared photon quantum states (polarization, frequency, phase). A photonic superposition
\begin{equation}
|\psi\rangle = \alpha|0\rangle + \beta|1\rangle
\end{equation}
can perform decomposition and recombination in quantum parallel, analogous to but richer than classical 0/1 bits \cite{Nielsen2010,Engel2007}.

\paragraph{(iii) Natural quantum communication.} When two subsystems are in pure quantum states, they can exchange information via coherence or entanglement without a classical channel \cite{Vedral2011}. If a trained human mind can sustain low-noise quantum coherence, it may couple informationally with macroscopic quantum systems such as the atmosphere or water.

\subsection{Mathematical description of information decomposition and recombination}

Let neural information $I$ in the coarse layer be represented by ionic and neurotransmitter fields $X(t)$. When mapped to the fine layer, $I$ decomposes into $N$ infrared modes:
\begin{equation}
|\Psi_I\rangle = \bigotimes_{k=1}^{N} \left(\alpha_k |0_k\rangle + \beta_k |1_k\rangle\right)
\end{equation}
where $|0_k\rangle$ and $|1_k\rangle$ denote two orthogonal quantum states of the $k$-th infrared mode. Processing is described by a unitary $U$:
\begin{equation}
|\Psi_O\rangle = U |\Psi_I\rangle
\end{equation}
and is then mapped back to the coarse layer as neural impulse or behavior. The physical realization of $U$ is governed by the coupling Hamiltonian between infrared photons and microtubule vibrational modes:
\begin{equation}
H = H_{\text{MT}} + H_{\text{IR}} + H_{\text{int}}
\end{equation}
with
\begin{align}
H_{\text{MT}} &= \sum_j \hbar\omega_j b_j^\dagger b_j, \\
H_{\text{IR}} &= \sum_k \hbar\nu_k a_k^\dagger a_k, \\
H_{\text{int}} &= \sum_{j,k} g_{jk} \left(a_k^\dagger b_j + a_k b_j^\dagger\right),
\end{align}
where $b_j^\dagger$, $b_j$ create/annihilate microtubule phonon/vibrational modes, $a_k^\dagger$, $a_k$ create/annihilate infrared photon modes, and $g_{jk}$ is the coupling strength \cite{Pokorny2004,Lambert2013}.

\subsection{Meditation, contemplation, and fine-quantum coherence}

During alert waking, neural noise (thermal fluctuations, ionic turbulence) is high and fine-quantum coherence decoheres rapidly. Meditation, contemplation, and inward stillness may reduce sympathetic arousal and synchronize electroencephalographic and respiratory rhythms, thereby lowering the decoherence rate $\Gamma_{\text{decoh}}$. We postulate a critical noise/temperature threshold $T_c$ below which the infrared coherence length $L_c$ increases sharply:
\begin{equation}
L_c(T) = v_g \tau_0 \exp\left(\frac{E_g}{k_B T_{\text{eff}}}\right)
\label{eq:coherence}
\end{equation}
where $v_g$ is the infrared waveguide group velocity, $\tau_0$ a characteristic scattering time, and $E_g$ an effective energy gap. This relation predicts stronger infrared coherence signatures with deeper meditative states.

\section{Cosmological Framework: Consciousness and the Cyclic Universe}

\subsection{The singularity as a consciousness node}

In standard cosmology the Big Bang singularity marks the breakdown of classical general relativity, with physical quantities diverging. Following Li's ``mathematical point, cosmic singularity, consciousness point'' framework \cite{Li2026a,Li2026b}, we propose that the singularity is not absolute nothingness but an information-bearing ``consciousness node.'' A consciousness field $\Psi_C$ is both outside the universe (transcending spacetime) and inside it (manifesting through biological consciousness). We write:
\begin{equation}
\Psi_C = \Psi_{\text{intra}} + \Psi_{\text{extra}}
\end{equation}
where $\Psi_{\text{intra}}$ is the component embedded in matter and $\Psi_{\text{extra}}$ transcends the classical spacetime boundary. The initial information of the singularity is given by the projection of $\Psi_C$ onto the boundary:
\begin{equation}
I_{\text{singularity}} = S\left(\rho_C|_{\partial M}\right)
\end{equation}
where $S$ is the von Neumann entropy and $\rho_C$ the reduced density matrix of the consciousness field.

\subsection{Mechanism of cosmic cycling}

Unlike Penrose's Conformal Cyclic Cosmology (CCC) \cite{Penrose2010,Gurzadyan2010}, our model emphasizes the driving role of the consciousness field:

\paragraph{(i)} The present universe evolves toward a ``Great Refinement'' state of high structural order and information density, with widespread biological consciousness.

\paragraph{(ii)} Information condenses into ``cosmic seeds,'' i.e., supermassive black holes that store information from the prior universe.

\paragraph{(iii)} A new universe is born from a bounce inside the black-hole horizon, carrying initial conditions imprinted by the preceding cycle and thereby avoiding heat death.

The cycle can be summarized as:
\begin{equation}
\text{Universe}_n \;\xrightarrow{\text{Great Refinement}}\; \text{Black-hole seed} \;\xrightarrow{\text{Bounce inside horizon}}\; \text{Universe}_{n+1}
\end{equation}

\subsection{``One thought, one universe'' and the Great Compassion Mantra}

We interpret ``one thought, one universe'' as a metaphor for feedback from the consciousness field to physical reality: when the brain changes the local configuration of $\Psi_C$ through the fine-quantum layer, that change may correlate with macroscopic events via quantum-holographic mechanisms. The Great Compassion Mantra and similar contemplative traditions are treated here not as physical equations but as cultural practices that may optimize fine-quantum coherence and resonance with the consciousness field.

\section{Methods, Simulations, and Testable Predictions}

\subsection{Numerical simulation protocol}

We propose a three-tier simulation strategy to test IRQIC:

\paragraph{(i) Microscopic scale.} Molecular dynamics (GROMACS/NAMD) of tubulin at 300 K to identify infrared-active vibrational modes in the 6--12 $\mu$m band; density-functional theory (DFT/VASP) to estimate coupling constants $g_{jk}$.

\paragraph{(ii) Mesoscopic scale.} Quantum-optical simulations (QuTiP or Strawberry Fields) of the Hamiltonian above to compute coherence time $\tau_\phi$, quantum fidelity $F$, and information throughput $C$.

\paragraph{(iii) Macroscopic scale.} Simultaneous EEG and infrared thermography to detect coherent infrared fluctuations during meditation or task states.

\subsection{Synthetic data example}

Using Eq.~(\ref{eq:coherence}), we generate a synthetic theoretical dataset (full code on GitHub). For room temperature 300 K, effective gap $E_g = 25$ meV, $\tau_0 = 10$ fs, and $v_g = c/4$, the infrared coherence length $L_c$ as a function of effective neural temperature $T_{\text{eff}}$ is shown in Table~\ref{tab:coherence}.

\begin{table}[htbp]
\centering
\caption{Synthetic theoretical coherence length $L_c$ versus effective neural temperature $T_{\text{eff}}$.}
\label{tab:coherence}
\begin{tabular}{@{}cc@{}}
\toprule
$T_{\text{eff}}$ (K) & $L_c$ (nm) \\ \midrule
250 & 18.2 \\
260 & 14.5 \\
270 & 11.7 \\
280 & 9.5  \\
290 & 7.8  \\
300 & 6.4  \\
310 & 5.3  \\ \bottomrule
\end{tabular}
\end{table}

The data show that lowering neural temperature/noise exponentially extends coherence length, providing a quantitative prediction for ``meditation-enhanced quantum coherence.''

\subsection{Testable predictions}

\paragraph{Prediction P1 (neural infrared superradiance).} Neural microtubules emit coherent infrared radiation in the 8--12 $\mu$m band within 1--10 ms after an action potential, with intensity correlated to spike frequency.

\paragraph{Prediction P2 (meditation-coherence enhancement).} Long-term meditators exhibit a significantly higher infrared thermography coherence index (IC-index) between prefrontal and parietal cortices than controls ($p < 0.05$).

\paragraph{Prediction P3 (black-hole information retention).} If the consciousness field participates in cosmic recycling, the end state of Hawking radiation should retain non-random patterns of information from the previous universe, testable by future gravitational-wave or cosmic-microwave-background anomalies.

\paragraph{Prediction P4 (non-zero singularity information).} If the singularity can carry information, the primordial power spectrum should differ from standard CCC predictions at specific scales, distinguishable by CMB B-mode polarization measurements.

\subsection{Benchmark comparison}

We benchmark IRQIC against three leading consciousness theories:

\paragraph{(i) Integrated Information Theory (IIT).} IIT stresses integrated information $\Phi$ \cite{Tononi2004}; IRQIC supplies a physical carrier for $\Phi$---namely infrared photon coherence---and explains how information integration may be physically implemented.

\paragraph{(ii) Global Workspace Theory (GWT).} GWT emphasizes prefrontal broadcasting \cite{Dehaene2011}; IRQIC adds a quantum broadcast channel beneath the classical workspace.

\paragraph{(iii) Orchestrated Objective Reduction (Orch-OR).} Orch-OR centers on microtubule quantum-gravitational reduction \cite{Hameroff2014}; IRQIC retains microtubules but replaces gravitational collapse with an infrared-photon channel, lowering the energy threshold and making the model more experimentally accessible.

\section{Data and Code Availability}

All simulation code, synthetic datasets, and figure-generation scripts are available at:
\begin{center}
\url{https://github.com/xiaohui-li-lab/2026-11-23/}
\end{center}

Core scripts:
\begin{itemize}
    \item \texttt{generate\_coherence\_data.py}: generates $L_c$ versus $T_{\text{eff}}$ data.
    \item \texttt{plot\_fig1\_dual\_scale.py}: schematic of the coarse/fine dual-scale model.
    \item \texttt{plot\_fig2\_irqic.py}: IRQIC model diagram.
    \item \texttt{plot\_fig3\_cyclic\_universe.py}: cyclic-universe schematic.
    \item \texttt{plot\_fig4\_coherence\_vs\_temperature.py}: Figure 4 curve.
\end{itemize}

\section{Discussion}

We have presented a cross-scale unifying hypothesis: consciousness processes neural information through an infrared photon fine-quantum layer at the microscale and drives a cyclic universe through a consciousness field at the cosmological scale. The theory bridges neuroscience, quantum optics, and cosmology, but it also faces major challenges: how to detect individual infrared photons carrying neural signals in vivo; how to distinguish quantum effects from classical thermal noise; and how to formalize the consciousness field within quantum gravity. Future progress depends on advances in experimental technologies (superconducting single-photon detectors, quantum imaging) and theoretical formalization (quantum field theory, holographic principles).

\section{Author Contributions}

\textbf{Xiaohui Li:} conceived the core concepts of ``infrared photons as the minimal units of conscious information,'' ``coarse/fine dual-scale processing,'' and ``consciousness field and cyclic universe''; wrote the manuscript; constructed the theoretical model and numerical simulation protocols; established the GitHub data repository.

\section{Competing Interests}

The author declares no competing interests.

\section{Ethics Statement}

This is a theoretical study and does not involve human or animal experiments.

\section{Acknowledgements}

The author thanks the open-science community for computational tools and literature resources.

\bibliographystyle{naturemag-doi}
\begin{thebibliography}{20}

\bibitem[Hodgkin \& Huxley(1952)]{Hodgkin1952}
Hodgkin A L, Huxley A F. A quantitative description of membrane current and its application to conduction and excitation in nerve. \textit{J Physiol}, 1952, 117(4): 500--544.

\bibitem[Waxman(1980)]{Waxman1980}
Waxman S G. Determinants of conduction velocity in myelinated nerve fibers. \textit{Muscle Nerve}, 1980, 3(2): 141--150.

\bibitem[Katz(1969)]{Katz1969}
Katz B. \textit{The release of neural transmitter substances}. Liverpool University Press, 1969.

\bibitem[Chalmers(1995)]{Chalmers1995}
Chalmers D J. Facing up to the problem of consciousness. \textit{J Conscious Stud}, 1995, 2(3): 200--219.

\bibitem[Hameroff \& Penrose(2014)]{Hameroff2014}
Hameroff S, Penrose R. Consciousness in the universe: a review of the `Orch OR' theory. \textit{Phys Life Rev}, 2014, 11(1): 39--78.

\bibitem[Craddock et al.(2018)]{Craddock2018}
Craddock T J A, Friesen D, Mane J, et al. A review of transcranial magnetic stimulation and the Orch-OR theory. \textit{Front Mol Neurosci}, 2018, 11: 433.

\bibitem[Popp et al.(1992)]{Popp1992}
Popp F A, Li K H, Gu Q. \textit{Recent advances in biophoton research and its applications}. World Scientific, 1992.

\bibitem[Fr\"{o}hlich(1968)]{Frohlich1968}
Fr\"{o}hlich H. Long-range coherence and energy storage in biological systems. \textit{Int J Quantum Chem}, 1968, 2(5): 641--649.

\bibitem[Li(2026a)]{Li2026a}
Li X.Infrared-Mediated Consciousness in Systems and the
Interoperability of Quantum-Induced Consciousness
Systems (2026-11-22, unpublished manuscript).

\bibitem[Pokorn\'{y}(2004)]{Pokorny2004}
Pokorn\'{y} J. Excitation of vibrations in microtubules in living cells. \textit{Bioelectrochemistry}, 2004, 63(1--2): 321--326.

\bibitem[Nielsen \& Chuang(2010)]{Nielsen2010}
Nielsen M A, Chuang I L. \textit{Quantum computation and quantum information}. Cambridge University Press, 2010.

\bibitem[Engel et al.(2007)]{Engel2007}
Engel G S, Calhoun T R, Read E L, et al. Evidence for wavelike energy transfer through quantum coherence in photosynthetic systems. \textit{Nature}, 2007, 446(7137): 782--786.

\bibitem[Vedral(2011)]{Vedral2011}
Vedral V. Living in a quantum world. \textit{Sci Am}, 2011, 304(6): 38--43.

\bibitem[Lambert et al.(2013)]{Lambert2013}
Lambert N, Chen Y N, Cheng Y C, et al. Quantum biology. \textit{Nat Phys}, 2013, 9(1): 10--18.

\bibitem[Li(2026b)]{Li2026b}
Li X.The Mathematical Point, the Cosmic Singularity, and
the Consciousness Point:
Infinitesimal Points as the Framework of the Universe(2026-11-21, unpublished manuscript).

\bibitem[Penrose(2010)]{Penrose2010}
Penrose R. \textit{Cycles of time: an extraordinary new view of the universe}. Bodley Head, 2010.

\bibitem[Gurzadyan \& Penrose(2010)]{Gurzadyan2010}
Gurzadyan V G, Penrose R. Concentric circles in WMAP data may provide evidence of violent pre-Big-Bang activity. arXiv:1011.3706, 2010.

\bibitem[Tononi(2004)]{Tononi2004}
Tononi G. An information integration theory of consciousness. \textit{BMC Neurosci}, 2004, 5(1): 42.

\bibitem[Dehaene \& Changeux(2011)]{Dehaene2011}
Dehaene S, Changeux J P. Experimental and theoretical approaches to conscious processing. \textit{Neuron}, 2011, 70(2): 200--227.

\bibitem[Maldacena \& Susskind(2013)]{Maldacena2013}
Maldacena J, Susskind L. Cool horizons for entangled black holes. \textit{Fortschr Phys}, 2013, 61(9): 781--811.


\textbf{\textbf{Declarations}}

  
  
  \textbf{\textbf{Acknowledgements}}

The author thanks colleagues for discussions on quantum aggregates, consciousness, emergence, and dimensionality. 

 \textbf{\textbf{Funding Statement}}:This research did not receive any specific grant from funding agencies in the public,commercial,or not-for-profit sectors.

Currently, there are no sponsors. Or I haven't found any sponsors or organizations yet. However, if I can publish it, I will cover part of the costs myself. Or it will be easier to find sponsors. (For example, the German company E.ON. I have previously found international buyers for their wind turbines.)

 \textbf{\textbf{Competing Interests}}

The author declares no competing interests.

 \textbf{\textbf{Author Contributions}}

Xiaohui Li: Conceptualization, Methodology, Investigation, Formal analysis, Writing – Original Draft, Visualization, Writing – Review \& Editing.

AI assisted me in rewriting the original text into an article that meets the submission requirements of Journal of the American Mathematical Society. No content was added or omitted.

 The author is solely responsible for the work described in this manuscript and has approved the final version for publication.

 \textbf{\textbf{Clinical trial:}}not applicable.

 \textbf{\textbf{Ethics and Consent to Participate declarations}}: not applicable.

 Hi,

This is an original work by the author. There is only one author:Xiaohui Li.

The article accurately reflects the scientific results.

Xiaohui Li agreed to submit the manuscript.

This manuscript has not been published in any other journals.

The author accepts the established procedure for selecting and publishing manuscripts.

good

Xiaohui Li



 \textbf{\textbf{Clinical trial number}}: not applicable.
 
 \textbf{\textbf{Declaration of Competing Interests}}: 

The authors declare that they have no competing interests.


\textbf{\textbf{Consent to Publish declaration missing}}: 

not applicable.


\end{thebibliography}

\end{document}
